% !TeX program = xelatex
% 采用与论文相同且未修改的 CUMCM 类文件和字体。
% 工具型号、使用日期及主要输入按参赛队补充信息填写。
\documentclass[withoutpreface,bwprint]{commons/cumcmthesis}
\setlength{\emergencystretch}{2em}
\pagestyle{plain}
\raggedbottom
\clubpenalty=10000
\widowpenalty=10000
\displaywidowpenalty=10000
\newcolumntype{Y}{>{\setstretch{1.5}\centering\arraybackslash}X}
\newcommand{\field}[1]{\mbox{\textbf{#1：}}}
\newenvironment{promptblock}{\begin{minipage}[t]{\linewidth}\setlength{\parindent}{0pt}}{\end{minipage}}
\setlist[enumerate]{leftmargin=2.6em,labelsep=0.4em,itemsep=0.75em,parsep=0pt}
\title{AI 工具使用详情}
\pdfstringdefDisableCommands{\def\thinspace{}}
\hypersetup{pdftitle={AI工具使用详情},pdfsubject={A题AI工具使用说明},pdfauthor={},pdfcreator={XeLaTeX}}

\begin{document}
\maketitle
\begin{center}
  {\songti\zihao{-4} A题：药材烘干过程的热湿迁移与收缩建模}
\end{center}

本说明结合参赛队补充信息及项目记录，列明 AI 工具的具体用途、主要交互方式和输出处理情况。使用以赛题、附件及队员建模文档为基础，围绕模型表达、程序实现与结果核对中的具体问题展开；参赛队负责模型选择、方案取舍、结果解释和最终稿审核。

\section{所用 AI 工具名称、版本或型号}
工具及使用信息见\cref{tab:ai-tools}。
\begin{table}[H]
\centering
\caption{AI 工具与使用信息}
\label{tab:ai-tools}
\normalfont\normalsize
\setlength{\tabcolsep}{9pt}
\setlength{\extrarowheight}{4pt}
\begin{tabularx}{0.96\textwidth}{@{}>{\setstretch{1.5}\centering\arraybackslash}p{0.22\textwidth}Y@{}}
\toprule
\textbf{项目} & \multicolumn{1}{c@{}}{\textbf{填写内容}} \\[5pt]
\midrule
工具名称 & OpenAI Codex / ChatGPT \\[6pt]
版本或型号 & GPT6 \\[6pt]
使用日期 & 2026年9月10日--13日 \\[6pt]
使用形式 & 对话讨论、代码辅助实现与调试、LaTeX 排版 \\[6pt]
主要输入 & 赛题及附件、队员建模文档、报错信息等 \\[5pt]
\bottomrule
\end{tabularx}
\end{table}

\section{具体使用目的和环节}
\subsection{模型表述与条件讨论}
以队员建模文档中的假设、控制方程和求解安排为输入，就径向离散、表面对流边界及收缩坐标表达中的疑点进行讨论。AI 输出包括条件解释、局部推导和实现建议，用于对照题目条件、检查公式与代码的对应关系。涉及收缩描述时，重点检查干基含水率、材料速度与网格速度的定义。对应论文第5节及第6至9节的模型部分。

\subsection{代码辅助实现与调试}
输入已有程序、公式、报错信息和运行结果，辅助处理附件读取、插值、矩阵组装、非线性迭代及首达时刻定位等实现问题，并辅助编写用于交叉校验的计算代码。修改后的程序通过实际运行、步长与网格对照检查；输出采样间隔与内部积分步长分别设置。对应论文第5.5至5.6节、各问求解部分及附录B。

\subsection{一致性复查}
以当前计算结果、论文数值表和文字中的关键数值为输入，辅助检查单位、时间位置、取整精度和数据来源的一致性，定位历史结果残留或转录差异。复查结论以源数据及校验记录为依据。对应论文第6至9节的结果部分。

\subsection{LaTeX 排版}
根据已有正文及编译信息，辅助调整公式、表格、标题层级和交叉引用，修正排版错误，并整理本使用说明。相关处理以准确呈现既有模型与计算结果为目的，涉及实质结论的修改须回到模型和结果文件核对。

\section{主要提示方式与使用过程说明}
\subsection{提示方式}
围绕具体问题提供必要材料，明确需检查的公式、代码片段或结果位置，再根据运行反馈补充条件、比较修改前后的差异。以下示例为代表性任务的归纳表述，并非原始对话的逐字记录。

\subsection{典型交互示例}
\begin{enumerate}[label=\textbf{（\arabic*）}]
  \item \begin{promptblock}\textbf{程序调试。}“根据建模文档中的方程与边界条件，检查这段有限体积程序的系数组装和迭代更新，说明报错原因及修改位置。”

  对应处理：比较建议与原公式，修订相关代码，再运行并检查数值范围、收敛情况和输出格式。
  \end{promptblock}

  \item \begin{promptblock}\textbf{收缩条件检查。}“检查已有收缩方程中材料速度、网格速度和干基含水率的定义是否一致，给出明确前提下的校验思路。”

  对应处理：补充速度与守恒关系的条件说明，采用配对计算检查实现一致性，并保留模型适用范围。
  \end{promptblock}

  \item \begin{promptblock}\textbf{一致性复查。}“逐项核对论文表格、正文关键数值和当前计算文件，列出不一致的位置、来源及建议修正值。”

  对应处理：回查源数据，确认差异属于计算设置、取整还是文字转录，再按当前结果统一展示值。
  \end{promptblock}
\end{enumerate}

\subsection{使用过程中的取舍}
建议的采纳依据是题目约束、模型含义和实际计算结果。代码修订结合运行反馈处理；改变边界、方程或结果解释的建议，连同其前提和对照结果讨论。方案取舍与物理解释以队员建模文档和实际计算证据为依据。

\section{对 AI 输出的采纳、人工修改和校验的主要情况}
本节列出语言润色之外的主要处理实例。项目记录可追溯相应修订和程序校验；人工复核按建模、编程及论文整理环节分工，重点确认建议与原模型的关系、修改理由及结论影响。

\subsection{积分与输出一致性：经修订采纳}
针对第二、三问共同时间段的结果检查，程序统一使用同一条积分轨迹，第三问按规定时刻抽取；终算采用161个径向节点和1秒步长。现有记录显示，共同时间层的温度、水分与阈值剖面一致，空间和时间加密检查通过。编程环节负责核对实现和输出，建模环节复核精度要求。该处理统一了离散与采样口径，首达判据保持为全域最大水分达到阈值。

\subsection{收缩条件与交叉校验：结合假设采纳}
收缩方程及验证实现的建议结合干基水分定义进行修订，明确均匀材料收缩、材料与网格同速等条件。项目保留均匀场、方程收支及移动控制体配对计算记录，用于检查实现的一致性。人工复核重点为守恒推导、变量含义和适用条件。结果解释采用相同物性下的固定半径对照，不将第三、四问时长之差全部归因于收缩，也不将数值吻合视为实验验证。

\subsection{正文与结果文件核对：核实后修正}
根据当前结果文件修正历史展示值，例如第三问6小时中心含水率由1.0156改为1.0170（$\mathrm{kg/kg}$）。2026年9月12日的专项记录对14张数值表和四份结果工作簿进行了核对，未发现与当前源值不符的项目。论文整理环节负责核对来源、单位和取整，编程环节复核数据对应关系；该次修正未重算仿真场，未改变既有首达判据。

\end{document}

